% SLLAB technical write-up. Build with: tectonic sllab.tex
% Layout follows Pascal Michaillat's latex-paper template (see sllab.sty).
\documentclass[letterpaper,11pt]{article}
\usepackage{sllab}
\bibliographystyle{sllab}

\hypersetup{pdftitle={SLLAB: A Platform for Running a College Club Sports Team}}

\newcommand{\code}[1]{\texttt{#1}}

\begin{document}

\title{SLLAB: A Platform for Running\\a College Club Sports Team}

\author{Lucius Scala%
\thanks{University of California, San Diego; \href{mailto:lscala@ucsd.edu}{lscala@ucsd.edu}.
\textit{AI disclosure}: parts of this codebase and this write-up were written with the help of AI.
All files were reviewed by a human, but they are provided as-is, with no guarantee that they are
correct, secure, or fit for any particular purpose.}}

\date{November 2026}

\paperurl{https://github.com/luciusscala/sllab}

\begin{titlepage}
\maketitle

College club sports teams are run by students, with no staff, little budget, and officers who
change every year or two. SLLAB is a platform that gives such a club the logistics of a staffed
program on a single backend. It consists of Fixture by SLLAB, an iOS app for schedules,
availability, reminders and attendance; UCSDMCS by SLLAB, a public club website with live
standings, a calendar feed, a web version of Fixture and an administration dashboard; and Film by
SLLAB, an iOS app that records games from a phone mounted on a custom motorized base, steered by
on-device object detection. The platform is in daily use by the UC San Diego Men's Club Soccer
team, and Fixture is available on the App Store. This write-up describes the system's
architecture and data model, the computer-vision and embedded control loop behind the camera,
the engineering practices used to maintain it, and the plan for generalizing one team's tools
into a platform that any club can adopt.

\end{titlepage}

% ============================================================================
\section{Introduction}\label{s:introduction}

\paragraph{The problem} A college club team does much of what a varsity program does: it
schedules practices, books fields, travels to tournaments, recruits at tryouts and keeps its
roster eligible. But students do all of it, in their spare time, with general-purpose tools:
group chats for attendance, spreadsheets for rosters, a form for tryouts, and a website that is
rarely updated. Each tool works on its own and nothing connects them, so the same information is
entered in several places and drifts out of date. When officers graduate, much of what they knew
leaves with them.

\paragraph{The cost of film} Game film is one of the most useful tools a team has for improving,
and one of the least accessible. It requires either a person behind the camera for every match or
a commercial auto-tracking camera, which costs around \$2{,}000.

\paragraph{The approach} SLLAB is one connected platform with one backend. The schedule a captain
creates is the schedule players respond to, the one the public website shows, and the one game
film is attached to. The design goals are that a club can adopt it without a developer, that it
works with the habits players already have, and that it is polished enough for the team to feel
like a program.

\paragraph{What exists} Four components are built. Fixture by SLLAB is an iOS app for
schedules, availability, reminders and attendance, released on the App Store at version~1.2.
UCSDMCS by SLLAB is the club's public website, which also hosts a web version of Fixture and an
administration dashboard. Film by SLLAB records games on a phone mounted on a motorized base and
is a working prototype. Hardware by SLLAB is that base: ESP32 firmware, a custom printed circuit
board and 3D-printed parts, costing about \$400 in total.

\paragraph{Outline} Section~\ref{s:overview} gives an overview of the system.
Sections~\ref{s:fixture} and~\ref{s:site} describe Fixture and the club website.
Section~\ref{s:data} describes the data model. Section~\ref{s:film} describes Film, its control
loop and its hardware. Section~\ref{s:practice} covers engineering practice, and
Section~\ref{s:roadmap} the roadmap. Appendix~\ref{a:timeline} gives a development timeline.

% ============================================================================
\section{System Overview}\label{s:overview}

Every component shares one Postgres database, hosted on Supabase \citep{Supabase}. The iOS apps
read and write the database directly, under access rules enforced by the database itself.
Operations that require a secret, such as sending text messages, run on a server instead.
Figure~\ref{f:architecture} shows the system as it is today, and Table~\ref{t:components} lists
each component's technology.

\begin{figure}[t]
\begin{tikzpicture}
  \def\xa{0}\def\xb{3.25}\def\xc{7.6}\def\xd{10.85}
  \node[box] (film) at (\xa,0) {\textbf{Film}\\iOS};
  \node[box] (fix)  at (\xb,0) {\textbf{Fixture}\\iOS};
  \node[box] (site) at (\xc,0) {\textbf{UCSDMCS}\\site, team app, admin};
  \node[plan] (api) at (\xd,0) {\textbf{API}\\planned};
  \node[box] (mount)  at (\xa,2.2) {\textbf{Mount}\\ESP32, stepper};
  \node[box] (yt)     at (\xb,2.2) {\textbf{YouTube}\\uploads};
  \node[box] (twilio) at (\xc,2.2) {\textbf{Twilio}\\SMS};
  \node[box] (se)     at (\xd,2.2) {\textbf{SportsEngine}\\standings};
  \node[db, minimum width=133.5mm] (db) at ({(\xa+\xd)/2},-2.2)
    {\textbf{Supabase}: Postgres, access rules, storage};
  \draw[arr] (film) -- node[note, left]{BLE} (mount);
  \draw[arr] (film.north east) -- (yt.south west);
  \draw[arr] (fix) -- node[note, above]{unlock} node[note, below]{reminders} (site);
  \draw[arr] (site) -- (twilio);
  \draw[arr] (se.south west) -- (site.north east);
  \draw[parr, <->] (film.south) -- (film.south |- db.north);
  \draw[arr, <->]  (fix.south)  -- (fix.south |- db.north);
  \draw[arr, <->]  (site.south) -- (site.south |- db.north);
  \draw[parr, <->] (api.south)  -- (api.south |- db.north);
\end{tikzpicture}
\caption{Architecture of the current system}\label{f:architecture}
\note{Solid lines are in use; dashed lines and boxes are planned. Fixture and the website share
the database. Fixture calls the website for the two operations that need a server: unlocking
captain mode and sending reminders. Film does not yet use the shared backend.}
\end{figure}

\begin{table}[t]
\caption{Components and their technology}\label{t:components}
\footnotesize
\begin{tabularx}{\linewidth}{@{}l X X@{}}
\toprule
Component & Technology & Role \\
\midrule
Fixture by SLLAB & SwiftUI, supabase-swift & Player and captain app \\
UCSDMCS by SLLAB & Next.js 16, React 19, Tailwind CSS 4 & Public site, web team app, administration, server endpoints \\
Film by SLLAB & SwiftUI, AVFoundation, Vision, Core ML, Core Bluetooth & Recording and tracking \\
Hardware by SLLAB & ESP32 (Arduino), TMC2209 driver, NEMA 17 stepper & Pan base and its controller \\
Detection model & Ultralytics YOLO26s, exported to Core ML & On-device detection \\
Database & Supabase, Postgres 17 & Single source of truth \\
\bottomrule
\end{tabularx}
\end{table}

% ============================================================================
\section{Fixture by SLLAB}\label{s:fixture}

Fixture is the app the team uses every day. Its central design constraint was adoption: a club
cannot ask thirty students to create accounts before they can say whether they are coming to
practice, so joining had to take seconds. Figure~\ref{f:fixture} shows its main screens.

\begin{figure}[t]
\subcaptionbox{Schedule}{\media{fixture-schedule.png}{0.225\linewidth}{58mm}{Season schedule with availability}}\hfill
\subcaptionbox{Event}{\media{fixture-event.png}{0.225\linewidth}{58mm}{An event with team responses}}\hfill
\subcaptionbox{Reminder}{\media{fixture-reminder.png}{0.225\linewidth}{58mm}{SMS reminder composer}}\hfill
\subcaptionbox{Attendance}{\media{fixture-dashboard.png}{0.225\linewidth}{58mm}{Captain attendance dashboard}}
\caption{Fixture by SLLAB on iOS}\label{f:fixture}
\end{figure}

\subsection{Players}

A player joins with a team code from their captain, picks and confirms their name on the roster,
and adds a phone number. From then on, the app opens to the season's schedule of practices and
games, each with its field, a link to directions and the player's own response. Marking oneself
going or not going takes one tap, and each event shows who else is coming.

\subsection{Captains}

Captain mode is unlocked with a separate code, which the server exchanges for a signed token. The
app keeps the token in the iOS Keychain and sends it with requests the server performs on a
captain's behalf, such as sending reminders. Captains can create games and practices, including
a repeating practice schedule in which choosing weekdays and a date range creates every session
at once. They can edit and delete events, see each event's responses grouped into going, not
going and no response, send reminders, and track attendance over the season.

\paragraph{Reminders never expose phone numbers} The app sends the server a list of roster
identifiers and a message, never phone numbers. The server looks the numbers up, sends the texts
through a Twilio Messaging Service, which also handles opt-out replies \citep{Twilio}, and
returns how many were delivered. Numbers never reach the captain's device, and no one outside the
roster can be texted through the endpoint.

\paragraph{Attendance is ranked by what a captain must act on} For a player with $n$ events in
the selected period, of which they declined $d$ and left $u$ unanswered, the dashboard orders
players by
\begin{equation}
\text{issue score} = 1.5\,\frac{u}{n} + \frac{d}{n}.
\label{e:issue}\end{equation}
An unanswered event weighs more than a declined one because a refusal lets a captain plan, while
silence does not. Each player's page lists the events they attended and those they missed, with
the reason.

\subsection{Operating concerns}

\paragraph{Universal links} Each reminder text contains a link of the form
\nolinkurl{app.ucsdmcs.com/events/<id>}, which opens the event directly in the app, or in the web
version if the app is not installed.

\paragraph{Mandatory updates} The database stores a minimum supported version, and older builds
show an update screen. The check fails open: if the app cannot reach the server, players are not
locked out.

\paragraph{Tests} Unit tests cover the models, decoding of the database's JSON, availability
grouping, attendance statistics and the data service.

% ============================================================================
\section{UCSDMCS by SLLAB}\label{s:site}

The club website serves three audiences from one Next.js deployment: the public, players using
Fixture on the web, and officers maintaining the season's data (Figure~\ref{f:site}).

\begin{figure}[H]
\subcaptionbox{Schedule and scores, desktop}{\media{site-schedule.png}{0.63\linewidth}{56mm}{Schedule and scores}}\hfill
\subcaptionbox{Standings, mobile}{\media{site-mobile.png}{0.3\linewidth}{56mm}{Standings}}
\caption{UCSDMCS by SLLAB}\label{f:site}
\note{The site is live at \url{https://www.ucsdmcs.com}.}
\end{figure}

\subsection{The public site}

The schedule lists each game with the opponent's crest, home or away, the field and the result,
beneath a season summary of record, win percentage, streak, home and away records and goal
difference. Standings come from the league's SportsEngine service. The site also carries the
roster, tryout information with a registration QR code, and the club's description. The schedule
is published as an iCalendar feed \citep{RFC5545}, so a game added in the dashboard appears in
subscribers' calendars without anyone re-sharing it.

\paragraph{Caching follows how the data changes} Public pages are cached without expiry and
invalidated by a single tag whenever the dashboard writes, because club data stays correct until
someone changes it. Standings are the exception: league administrators enter results on their
own schedule, so standings are refreshed at most once a day.

\subsection{The team app and the dashboard}

The web team app is a port of Fixture served at \code{app.ucsdmcs.com}. A single deployment
serves both the club site and the app: a proxy step reads the request's host and serves the team
section from the root, so each has its own address without a separate deployment.

The web and iOS clients share one session format, a payload signed with HMAC-SHA256
\citep{RFC2104}. The website stores it in an HTTP-only cookie and the app sends it as a bearer
token, so one format and one verification path serve both.

The dashboard lets officers edit teams and crests, the roster, the schedule, practices, socials,
tournaments and fields. Every write is a server action that checks the session itself rather than
relying on the page layout, because server actions can be invoked directly.

% ============================================================================
\section{Data Model}\label{s:data}

The schema separates durable identity from seasonal participation, and generic events from
their type-specific details. Figure~\ref{f:schema} shows the core tables.

\begin{figure}[t]
\begin{tikzpicture}
  \NewDocumentCommand{\tb}{m m m m}{%
    \node[tbody, anchor=north] (#1-b) at #2 {#4};
    \node[thead, anchor=south] (#1-h) at (#1-b.north) {\strut#3};
    \node[fit=(#1-h)(#1-b), inner sep=0pt] (#1) {};
  }
  \tb{teams}{(0,0)}{teams}{org, sport\\name, division\\crest, team\_code}
  \tb{seasons}{(4.4,0)}{seasons}{year\\is\_current}
  \tb{people}{(8.8,0)}{people}{name, hometown\\phone, picture}
  \tb{orgs}{(0,-2.7)}{orgs}{name}
  \tb{events}{(4.4,-2.7)}{events}{team, season\\type, date}
  \tb{roster}{(8.8,-2.7)}{roster}{person, team, season\\number, position\\class}
  \tb{sub}{(0,-5.3)}{event subtypes}{games: home, away,\\\quad score\\practices: notes\\tournaments, socials}
  \tb{avail}{(8.8,-5.3)}{availability}{event, roster\\status}
  \tb{fields}{(0,-8.3)}{fields}{name, address\\parking, picture}
  \draw[arr] (teams.south) -- (orgs.north);
  \draw[arr] (events.north) -- (seasons.south);
  \draw[arr] (events.north west) -- (teams.south east);
  \draw[arr] (roster.north) -- (people.south);
  \draw[arr] (roster.north west) -- (seasons.south east);
  \draw[arr] (roster.east) -- ++(0.45,0) |- ([yshift=4mm]teams.north) -- (teams.north);
  \draw[arr] (avail.north) -- (roster.south);
  \draw[arr] (avail.north west) -- (events.south east);
  \draw[arr] (sub.north east) -- (events.south west);
  \draw[arr] (sub.south) -- node[note, right]{field} (fields.north);
  \draw[arr] (sub.west) -- ++(-0.45,0) |- node[note, pos=0.25, left]{home,\\away} (teams.west);
\end{tikzpicture}
\caption{Core tables of the schema}\label{f:schema}
\note{Arrows point from a row to the row it references. Games, practices, tournaments and socials
each have their own table keyed by the event; they are grouped here.}
\end{figure}

\paragraph{People are not roster entries} A person exists once. A roster entry is that person on
one team in one season, with that season's number, position and class. A player's history, and
later their career statistics and alumni status, follow from this separation.

\paragraph{Events are a supertype} Every event has a team, a season, a type and a time, so the
schedule, availability and reminders work identically for all of them. Details specific to games
or practices live in subtype tables keyed by the event.

\paragraph{Availability belongs to a roster entry and an event} A uniqueness constraint on the
pair makes each response an upsert, so responses cannot be duplicated.

\paragraph{Games name two teams} A game records a home and an away team rather than ``us'' and an
opponent. It reads as home or away depending on whose side it is viewed from, so the same row can
serve both clubs once both are on the platform.

\paragraph{Constraints live in the database} Allowed positions, classes and statuses are check
constraints, and a partial unique index allows only one current season.

\paragraph{Limitations} Several of these choices were right for one team and must change for
many. The current season is global rather than per team, positions are those of soccer, and
joining relies on shared team and captain codes rather than personal accounts.
Section~\ref{s:roadmap} describes how each changes.

% ============================================================================
\section{Film by SLLAB and the Tracking Mount}\label{s:film}

Film turns a phone into an unattended camera operator. The phone records on its ultra-wide
camera while a detection model locates play in the frame, and a motorized base turns to keep it
centered (Figure~\ref{f:hardware}).

\begin{figure}[H]
\subcaptionbox{Mount with phone}{\media{mount.jpg}{0.31\linewidth}{48mm}{The assembled mount}}\hfill
\subcaptionbox{Controller board}{\media{pcb.jpg}{0.31\linewidth}{48mm}{The assembled PCB}}\hfill
\subcaptionbox{Tracking test}{\media{tracking.png}{0.31\linewidth}{48mm}{A frame from a tracking test}}
\caption{Hardware by SLLAB}\label{f:hardware}
\end{figure}

\subsection{The control loop}

Figure~\ref{f:loop} shows the path from camera frame to motor.

\begin{figure}[H]
\begin{tikzpicture}[node distance=3mm]
  \tikzset{box/.append style={minimum width=0mm, inner xsep=3pt, font=\scriptsize}}
  \node[box] (cam) {Camera\\every frame};
  \node[box, right=of cam] (det) {YOLO26s\\every 3rd frame};
  \node[box, right=of det] (sel) {Target\\best person};
  \node[box, right=of sel] (ema) {Smoothing\\EMA};
  \node[box, right=11mm of ema] (esp) {ESP32\\deadzone};
  \node[box, right=of esp] (mot) {Stepper\\3.4° step};
  \draw[arr] (cam) -- (det); \draw[arr] (det) -- (sel); \draw[arr] (sel) -- (ema);
  \draw[arr] (ema) -- node[note, above]{1 byte} node[note, below]{BLE} (esp);
  \draw[arr] (esp) -- (mot);
  \begin{scope}[on background layer]
    \node[fit=(cam)(ema), fill=gray!8, inner sep=4pt, label={[note]below:iPhone}] {};
    \node[fit=(esp)(mot), fill=gray!8, inner sep=4pt, label={[note]below:mount}] {};
  \end{scope}
\end{tikzpicture}
\caption{The tracking control loop}\label{f:loop}
\end{figure}

\paragraph{Detection} Frames are passed to a YOLO26s model \citep{Ultralytics} through Apple's
Vision and Core ML frameworks. Only every third frame is processed, which keeps inference within
budget while the phone also records. The target is the most confident person detection above a
confidence of 0.4.

\paragraph{Error signal} Let $x_c \in [0,1]$ be the target's horizontal center in normalized
image coordinates. The offset from the center of the frame is
\begin{equation}
o_t = 2\,(x_c - 0.5) \in [-1, 1],
\label{e:offset}\end{equation}
and it is smoothed with an exponential moving average, so that a single noisy detection does not
swing the camera:
\begin{equation}
s_t = 0.7\,s_{t-1} + 0.3\,o_t.
\label{e:ema}\end{equation}
When nothing is detected, the average resets and the base holds still until play is found again.

\paragraph{Protocol} The smoothed offset is quantized to one signed byte,
$b = \operatorname{trunc}(127\,s_t)$, and written over Bluetooth Low Energy without
acknowledgement. A single byte with no round trip keeps latency low, and a lost command is
superseded by the next one a fraction of a second later.

\paragraph{Firmware} The ESP32 keeps only the newest command. Commands that arrive while the motor
is moving overwrite one another instead of queueing, so the mount never moves toward where the
target \emph{was}. Offsets with $|b| \le 15$, about 12\% of full scale, are ignored, so the
camera holds steady when play is roughly centered instead of hunting back and forth. Outside this
deadzone, the motor moves 15 microsteps toward the target. With the driver at $1/8$
microstepping on a 200-step motor, each move turns the camera by
\begin{equation}
\frac{15}{200 \times 8} \times 360^\circ \approx 3.4^\circ.
\label{e:step}\end{equation}

\subsection{Hardware}

The base is a NEMA~17 stepper motor driven by a TMC2209, whose quiet mode keeps motor noise off
the recording's audio. The controller is a custom two-layer, through-hole board of about
$45 \times 65$~mm, designed in EasyEDA, with the ESP32 and the driver in header sockets so either
can be replaced without soldering. A 12~V supply feeds both the motor and the ESP32, with a
100~µF bulk capacitor to absorb current spikes when the motor steps. The housing, shaft coupler
and phone holder are 3D-printed. The design progressed from a hand-wired breadboard to a printed
base and then to the custom board, and the enclosure and board were refined with a mechanical
engineer and a systems engineer. The complete system costs about \$400.

\subsection{Recording, library and upload}

Film records immediately or on a timer. Videos and thumbnails are kept in the app's own library
rather than the camera roll and are excluded from iCloud backup, with metadata stored in
SwiftData. Videos are uploaded to YouTube as unlisted through the YouTube Data API's resumable
upload protocol \citep{YouTubeUpload}, and the resulting link is saved with the video.

\subsection{The detection model}

The shipping model is YOLO26s, tracking people. To track the ball instead, the model is being
fine-tuned with Ultralytics on a public soccer-ball dataset \citep{RoboflowBall}, at
$640 \times 640$ for 100 epochs on the Apple GPU. A second dataset, which labels the ball,
goalkeepers, players and referees \citep{RoboflowPlayers}, is prepared for later work. The
trained model is exported to Core ML with non-maximum suppression built in, so the app receives
final detections.

\subsection{Before release}

Film is a working prototype, and four changes precede its release. Upload authorization moves to
the platform's API, so that each team uploads to its own channel and the app holds no
credentials. Uploads stream from disk instead of loading whole files into memory, which
full-length matches require. Tracking moves from the most confident player to the ball. Finally,
each video is linked to its game through the shared backend, so that film appears in Fixture.

% ============================================================================
\section{Engineering Practice}\label{s:practice}

\paragraph{One repository} The project began as five repositories, which were merged into one
with their full history. A schema change and every client it affects can now land in a single
commit. Each application keeps its own folder and toolchain, and code moves into shared packages
only when a second consumer needs it.

\paragraph{The schema is code} The live database was captured as a baseline migration. Every
later change is an ordered SQL file applied with the Supabase command-line tool, never an edit in
the dashboard. The tool's version is pinned, because a major release changed the default target
of database commands from the hosted project to the local one.

\paragraph{Secrets stay out of version control} Every repository's history was scanned for keys
before the merge, and only templates of required variables are committed. The module holding the
server's secret database key is marked server-only, so importing it from client code is a build
error rather than a convention.

\paragraph{Large artifacts stay out of version control} Training data, runs and weights are
excluded. The repository holds the scripts and dataset definitions needed to re-create them.

\paragraph{Continuous integration} A path-filtered workflow builds and lints each web application
when it changes.

% ============================================================================
\section{Roadmap}\label{s:roadmap}

The next stage turns one team's tools into a platform that any club can join without a
developer. Figure~\ref{f:target} shows the target structure, and the steps below are in order.

\begin{figure}[H]
\begin{tikzpicture}[node distance=3.5mm]
  \tikzset{every node/.append style={minimum width=23mm}}
  \node[box] (fix) {\textbf{Fixture}\\iOS};
  \node[box, right=of fix] (film) {\textbf{Film}\\iOS};
  \node[box, right=of film] (web) {\textbf{Platform web}\\console, editor};
  \node[box, right=of web] (sites) {\textbf{Club sites}\\every club};
  \node[box, right=of sites] (api) {\textbf{API}\\logic, integrations};
  \node[db, below=12mm of web, minimum width=129mm] (db)
    {\textbf{One database}: accounts, memberships, teams, events, people, matches, messages, sites};
  \foreach \n in {fix,film,web,sites,api} \draw[arr, <->] (\n.south) -- (\n.south |- db.north);
\end{tikzpicture}
\caption{Target structure of the platform}\label{f:target}
\note{Every client uses one backend, and platform logic lives in one API.}
\end{figure}

\begin{enumerate}
\item \textit{Accounts and access.} Verified sign-in by phone or email, and a membership table
relating each account to its teams and roles. Access rules in the database then scope every read
and write to the caller's teams, and captain becomes a role granted by another captain rather
than a shared code. Fixture's mandatory-update mechanism moves every player to the new version
before the old paths are removed.
\item \textit{A platform API.} Server operations move from the club's website into a dedicated
service with versioned routes, so no application depends on a club's domain. Messaging becomes a
subsystem with templates, audience segments, scheduled sends, consent and delivery records.
\item \textit{Multi-team data.} Seasons per team, sport-specific positions, and a shared directory
of schools, teams, crests, fields and leagues that every club references instead of copying.
Each new club then improves the data for the others.
\item \textit{Matches and statistics.} Lineups and match events, from which player profiles,
career statistics and standings follow.
\item \textit{Club sites as data.} One deployment renders every club's site from its domain, with
content stored as typed blocks and themed with the club's colors and crest. The UC San Diego site
becomes the first set of rows instead of a codebase.
\item \textit{Graphics, gear and dues.} Match-day and result graphics generated from match data;
kit orders that collect sizes against roster numbers; and dues collected through each club's own
payment account.
\item \textit{Film on the platform.} Videos attached to their events, from Film or any other
source, and visible in Fixture.
\end{enumerate}

% ============================================================================
\bibliography{sllab}

% ============================================================================
\section{Development Timeline}\label{a:timeline}

Table~\ref{t:timeline} summarizes development, from the projects' commit histories.

\begin{table}[H]
\caption{Development timeline, 2026}\label{t:timeline}
\footnotesize
\begin{tabularx}{\linewidth}{@{}l X@{}}
\toprule
Month & Milestone \\
\midrule
April & Film begins: camera capture and live preview. \\
June & Detection runs on device, from frame capture to bounding boxes. \\
July & The phone steers the mount over Bluetooth. \\
August & Recording moves to the ultra-wide camera, and tracking works at 10--15~ft. In-app library and YouTube upload. The sllab.io site launches. \\
September & The club website launches with the dashboard, live standings and tryouts. Fixture enters beta at version 1.0.0, followed by its web version. \\
October & Text-message reminders and Keychain-backed captain sessions. Attendance tracking for captains. A revised schema, unit tests and mandatory updates. \\
November & Past events in Fixture. The mount's firmware, board and CAD files are published. The five repositories are merged into one. \\
\bottomrule
\end{tabularx}
\end{table}

\end{document}
